\chapter{Related Work and State of the Art}
\label{chapter:soa&rw}

\section{Related Work}
\label{sec:rw}
At the moment of writing, there is no online evidence about any unikernel based on a game server. Instead, a few unikernel projects are briefly described, which
can be candidates for this job.
\subsection{MirageOS}
MirageOS\footnote{\url{https://mirage.io/}} is a library operating system written in the OCaml language that constructs unikernels for secure, high-performance networking applcations.
\subsection{IncludeOS}
IncludeOS\footnote{\url{https://www.includeos.org/}} allows you to run your application in the cloud without an operating system. It adds operating system functionality to your application
allowing you to create performant, secure and resource efficient virtual machines.
\subsection{OSv}
OSv\footnote{\url{https://osv.io/}} is an open-source modular unikernel designed to run unmodified Linux applications securely on micro-VMs in the cloud.
\section{State of the Art}
\label{sec:soa}
So far, the following steps have been completed:
\begin{itemize}
    \item Ran an OpenArena server from scratch (source files).
    \item Ported the server to Unikraft in binary compatibility mode using the existing libelf library and a Dockerfile-based root filesystem.
    \item Ported the server to Unikraft in native mode by creating an external library containing the required source files and dependencies.
\end{itemize}